%! Author = sriadityavedantam
%! Date = 3/30/26

\section{Congruences Revisited}

\begin{definition}[Congruences]
    Let $m(x)\in\ff[x]$ with $m(x)\neq 0$.
    If $a(x),b(x)\in\ff[x]$, we define
    \[
        a(x)\equiv b(x)\mod m(x)
    \]
    if
    \[
        m(x)\mid a(x)-b(x).
    \]
    Equivalently, there exists $q(x)\in\ff[x]$ such that
    \[
        a(x)-b(x)=q(x)m(x).
    \]
    Equivalently,
    \[
        a(x)=b(x)+q(x)m(x).
    \]
\end{definition}

\begin{definition}[Congruence Class]
    The congruence class of $a(x)\in\ff[x]$ is denoted by $[a(x)]$.
    It consists of
    \[
        [a(x)] \coloneqq \{b(x)\in\ff[x] : b(x)\equiv a(x)\mod m(x)\}.
    \]
\end{definition}

\begin{definition}[Polynomial Division Algorithm]
    Suppose $g(x)\in\ff[x]$.
    Then
    \[
        g(x)=q(x)m(x)+r(x),
    \]
    where $\deg r(x)<\deg m(x)$ or $r(x)=0$.
    Since
    \[
        g(x)-r(x)=q(x)m(x),
    \]
    we have
    \[
        r(x)\equiv g(x)\mod m(x).
    \]
    So $g(x)\in [r(x)]$.
    Thus every $g(x)$ is in exactly one of these congruence classes.
\end{definition}

\begin{lemma}{
    In $\z_p[x]$, if $\deg m(x)=n$, then there are exactly $p^n$ different congruence classes modulo $m(x)$.
}
    Every polynomial is congruent modulo $m(x)$ to exactly one polynomial of degree less than $n$ or to $0$.
    Such a representative has the form
    \[
        a_0+a_{1}x+\cdots+a_{n-1}x^{n-1},
    \]
    where each coefficient is in $\z_p$.
    There are $p$ choices for each coefficient, and there are $n$ coefficients.
    So there are exactly
    \[
        p^n
    \]
    such polynomials.
    Hence there are exactly $p^n$ distinct congruence classes.
\end{lemma}

Similar to congruence classes in the integers, we also have similar ideas for addition and multiplication for polynomial congruences.

\begin{definition}[Modular Operations]
    Addition is defined by
    \[
        [a(x)]+[b(x)]\coloneqq [a(x)+b(x)].
    \]
    Multiplication is defined by
    \[
        [a(x)][b(x)]\coloneqq [a(x)b(x)].
    \]
\end{definition}

We can use this to check that these operations are well-defined.

\begin{lemma}[Well-Defined]{
    Suppose $[a(x)]=[c(x)]$ and $[b(x)]=[d(x)]$.
    Then:
    \begin{enumerate}
        \item $[a(x)+b(x)]=[c(x)+d(x)]$;
        \item $[a(x)b(x)]=[c(x)d(x)]$.
    \end{enumerate}
}
    Since $[a(x)]=[c(x)]$, we have
    \[
        a(x)\equiv c(x)\mod m(x).
    \]
    Since $[b(x)]=[d(x)]$, we have
    \[
        b(x)\equiv d(x)\mod m(x).
    \]
    Thus
    \[
        m(x)\mid a(x)-c(x)
    \]
    and
    \[
        m(x)\mid b(x)-d(x).
    \]
    Therefore
    \[
        m(x)\mid (a(x)-c(x))+(b(x)-d(x))=(a(x)+b(x))-(c(x)+d(x)).
    \]
    So
    \[
        a(x)+b(x)\equiv c(x)+d(x)\mod m(x),
    \]
    which proves
    \[
        [a(x)+b(x)]=[c(x)+d(x)].
    \]
    Also,
    \begin{align*}
        a(x)b(x)-c(x)d(x)
        &= a(x)b(x)-c(x)b(x)+c(x)b(x)-c(x)d(x)\\
        &= (a(x)-c(x))b(x)+c(x)(b(x)-d(x)).
    \end{align*}
    Since $m(x)\mid a(x)-c(x)$ and $m(x)\mid b(x)-d(x)$, it follows that
    \[
        m(x)\mid a(x)b(x)-c(x)d(x).
    \]
    So
    \[
        a(x)b(x)\equiv c(x)d(x)\mod m(x),
    \]
    which proves
    \[
        [a(x)b(x)]=[c(x)d(x)].
    \]
\end{lemma}